\paragraph{Matrix gauge realisation of the compensating rotation, the superpotential of the stress and the edge term (conditional; not closure).}
Keep the setting of the preceding paragraph: four matrix coordinates, the Poisson limit, \(\omega=\Theta^{-1}\) closed, \(n_D=0\), the \(\LUV^2\) density, \(c_W(r)=N_{\mathrm{tot}}/6+\kappa_Xr\), \(\xi\) the rotation (boost) of the normal plane of \(\Sigma\), \(B^a{}_b=\partial_b\xi^a\), and \(\{f,g\}=\Theta^{\mu\nu}\partial_\mu f\,\partial_\nu g\).

\emph{Gauge realisation.} If \(\mathcal L_\xi\omega=0\) on a star-shaped domain, then \(\iota_\xi\omega\) is closed and \(\iota_\xi\omega=\mathrm d\varepsilon\) with \(\varepsilon=0\) on \(\Sigma\); there \(\mathrm d\varepsilon=0\) and the Hessian of \(\varepsilon\) is \(B^{T}\omega\neq0\), and \(\{\varepsilon,f\}=\xi^\mu\partial_\mu f\) for every function \(f\). In the matrix chart \(\{\varepsilon,X^a\}=B^a{}_bX^b\): the rigid rotation of the preceding paragraph is the Poisson-limit matrix gauge transformation generated by \(\varepsilon\), for every background invariant under the joint flow and in both signatures (Poincar\'e lemma; exact checks for \(\omega=\bar\omega+\mathrm dA\) with \(A\) a random polynomial invariant one-form). For constant \(\bar\omega\) this is the Poisson limit of the quadratic generator of the constant stabiliser in Section~\ref{sec:status-D}. It follows from the fact that in the Poisson limit gauge transformations act on the \(X^a\) as diffeomorphisms preserving \(\omega\); it adds no symmetry and is an isotropy of the background, not a local rotation of the matrix index.

\emph{Gauge-covariant charge.} With the \(X^a\) as fields on the leaf and \(\omega\) fixed, \(\delta_\varepsilon X^a=\mathcal L_\zeta X^a\) with \(\zeta\) the Hamiltonian field of \(\varepsilon\), so the Noether current and charge of \(\delta_\varepsilon\) are those of the diffeomorphism \(\zeta\), with the same ambiguities, and here \(\zeta=\xi\). The combination of \(\xi\) with the compensating transformation \(-\varepsilon\), the analogue of the construction of~\cite{JacobsonMohd2015}, leaves every field invariant and its charge vanishes identically. The charge of \(\delta_\varepsilon\) alone equals that of \(\xi\): on \(\Sigma\) it depends on the second derivatives of \(\varepsilon\) and is nonzero for a curvature density, although \(\varepsilon\) and \(\mathrm d\varepsilon\) vanish there. In the variables \(\Theta(x)\) of the matrix chart, which these transformations leave invariant, the remaining freedom is that of a potential \(A\) of \(\omega=\bar\omega+\mathrm dA\); the compensating parameter \(\iota_\xi A+\varepsilon-\bar\varepsilon\), \(\mathrm d\bar\varepsilon=\iota_\xi\bar\omega\), vanishes on \(\Sigma\), and so does its charge when \(A\) enters only through \(\omega\), leaving the charge of the preceding paragraph. In both readings the term \(Y(\phi,\mathcal L_\xi X)=Y(\phi,\delta_\varepsilon X)\) is unchanged and \(r\) is not selected.

\emph{Superpotential of the stress.} On a flat background the identically conserved symmetric tensors linear in \(h\) with two derivatives are the combinations of the linearised Einstein tensor and \((\partial_\mu\partial_\nu-\eta_{\mu\nu}\Box)h^\alpha{}_\alpha\), and the divergence \(\partial^\rho\Sigma_{\rho\mu\nu}\) of every \(\Sigma\) linear in \(\partial h\), antisymmetric in \((\rho,\mu)\) and with symmetric divergence, is such a combination (exact classification, both signatures). Each is \(\partial^\rho\partial^\sigma U_{\mu\rho\nu\sigma}\) with \(U\) linear in \(h\) and with the symmetries of a curvature tensor. For a Killing field \(\xi\) the current \(\xi^\nu\partial^\rho\partial^\sigma U_{\mu\rho\nu\sigma}\) is the divergence of an antisymmetric tensor that reduces on \(\Sigma\) to \(\partial^\rho\xi^\nu U_{\mu\rho\nu\sigma}\); for the conformal Killing field of a diamond a bulk term in \(\partial\partial\xi\) remains (exact checks). On traceless \(T_{\mu\nu}\) the superpotential term of Equation~\eqref{eq:Top-FP} is therefore \(\lambda\,\mathcal E_{\mu\nu}[T]\): its edge contribution changes the coefficient of \(\mathcal E\) and is not independent of it. For \(T^{\mathrm{op}}_{\mu\nu}\) defined as the response of the induced quadratic action to \(G\) on the two modes with \(\delta\rho=0\), the Euler operator annihilates divergences and the coefficient is \(c_T=c_W(-4)\), so the edge term carried by \(T^{\mathrm{op}}_{\mu\nu}\) so defined is that of \(r=-4\), which restates the definition of \(r=-4\) by agreement on those modes. With the area coefficient \(c_A=c_W(0)\) of the fixed-state replica, the requirement \(G_A=G_T\) of the paragraph on the normalisation then needs the further term \(-4\kappa_X\) from the modular Hamiltonian of the state; the stress, so defined, does not supply it.

\emph{Edge terms of the fluctuation Lagrangians.} For a thermal state the modular Hamiltonian is the integrated Noether current of the quantum fields plus a term at the outer boundary; relative to the integral of the stress tensor it differs on \(\Sigma\) by the Noether charge of the matter Lagrangian, with the ambiguities of that charge set to zero, and by a contribution of the background dependence of the path-integral measure~\cite{Chen2022}. The tree-level fluctuation Lagrangians of the Poisson-limit paragraphs depend on the background through \((\gamma,\omega)\) and their first derivatives only, so their Noether charge on \(\Sigma\) vanishes, by the argument of the preceding paragraph, for fluctuations regular there. The measure contribution is defined there through a deformation localised near \(\Sigma\) that multiplies the opening angle by \(n\); about constant \(\bar\Theta\), with \(\Theta\) the only field, a localised first-order perturbation changes no opening angle, because the integral of \(R_G^{(1)}\) over the normal plane is \(-2\) times the flux of \(\nabla\ln\rho\) (paragraph on replica prescriptions), so that construction is not available in this class. In the Moser variables the linearised gauge transformation \(\delta a^\mu=\Theta^{\mu\rho}\partial_\rho\lambda\) of Remark~\ref{rem:ch3-trace-quadratic} acts as \(\delta A=-\mathrm d\lambda\) on \(A_\nu=a^\mu\omega_{\mu\nu}\), and on the stationary set the connection operator is the Maxwell--axion operator of the paragraph on Moser variables, whose curvature coupling is the Ricci term of the Hodge operator of \(G\), so the \(U(1)\) structure and the Ricci coupling on which the Kabat value \(-1\) of the vector contact rests are present there. If the state carries that edge sector, the edge-mode term of~\cite{DonnellyWall2016} for a Maxwell field, the area coefficient becomes \(c_A-1=c_W(-1)\) and \(3\kappa_X\) remains: the part that the stationary-set identity of the paragraph on Moser variables assigns to the gradient of \(\vartheta\) when that gradient is kept as an independent invariant, a first-order invariant with no charge on \(\Sigma\) in that presentation; in the presentation that eliminates it by the identity this part is carried by the charge of \(R_G\) with \(r=-4\), and the density does not select between the two.

These statements are for the Poisson limit, four matrix coordinates, \(n_D=0\), the \(\LUV^2\) density, backgrounds invariant under the joint flow and regular on \(\Sigma\), local charges, linear order about a flat background for the classification, and the cited construction as stated; the checks are exact at the sampled backgrounds, and the stationary-set identity used is numerical. They do not determine the edge term of the matrix state, its modular Hamiltonian or the replica, do not treat Moyal corrections, under which gauge transformations with non-quadratic generators are not diffeomorphisms, more than four matrix coordinates, the Dirac sector or the measure of the matrix integral, and derive no local Lorentz symmetry. Problems~B and~D remain open and PROGRAMME-CONDITIONAL.
